\documentclass[11pt,a4paper]{article}

\usepackage[utf8]{inputenc}
\usepackage[T1]{fontenc}
\usepackage{lmodern}
% \usepackage{geometry}
\usepackage{amsmath,amssymb,bm}
\usepackage{enumitem}
\usepackage{hyperref}
\usepackage{verbatim}


\usepackage{microtype}
\usepackage[margin=2.5cm]{geometry}
\usepackage{xcolor}
\usepackage{booktabs}
\usepackage{array}
\newcolumntype{L}[1]{>{\raggedright\arraybackslash}p{#1}}

\setlength{\parskip}{0.4em}
\setlength{\parindent}{0pt}
\renewcommand{\labelenumi}{(\alph{enumi})}

\definecolor{ink}{gray}{0.10}
\definecolor{mute}{gray}{0.40}
\definecolor{hair}{gray}{0.75}
\definecolor{panel}{gray}{0.94}

\newcommand{\R}{\mathbb{R}}

% ---- plain light-gray note/context box ----
\newsavebox{\ctxboxsv}
\newenvironment{ctxbox}[1]{%
  \par\smallskip\noindent
  \begin{lrbox}{\ctxboxsv}
  \begin{minipage}{\dimexpr\linewidth-2\fboxsep-2\fboxrule\relax}
  \textbf{#1}\par\smallskip
}{%
  \end{minipage}
  \end{lrbox}%
  \noindent\colorbox{panel}{\usebox{\ctxboxsv}}\par\smallskip
}

\geometry{margin=2.5cm}

\title{\bfseries Lab Assignment 1\\[4pt]
\Large Vectors, Matrices, and Linear System Solving}
\author{Maths II -- Bachelor's Degree in Video Game Design and Development\\CITM - UPC}
\date{October 2026}

\begin{document}

\maketitle

\begin{ctxbox}{Abstract}
This assignment introduces a practical workflow for core linear algebra and numerical
methods. Students implement custom \texttt{Vector} and \texttt{Matrix} classes, basic vector-space
and matrix-space operations, and a linear-system solver based on Gaussian elimination,
back substitution, and residual analysis. The work is validated both mathematically and
computationally through an interactive app and automated tests. The goal is to connect
formal mathematical definitions with robust C++ implementations, including dimension
checks, tolerance handling, and pivoting strategies.
\end{ctxbox}

\section*{Exercise 1: Vector and Matrix Methods}

\subsection*{a) Mathematical Summary}
A vector in $\mathbb{R}^n$ supports:
\begin{itemize}[leftmargin=*]
    \item Addition: $(u+v)_i = u_i + v_i$
    \item Subtraction: $(u-v)_i = u_i - v_i$
    \item Scalar product by $\lambda \in \mathbb{R}$: $(\lambda u)_i = \lambda u_i$
    \item Dot product: $u \cdot v = \sum_{i=0}^{n-1} u_i v_i$
\end{itemize}

A matrix in $\mathbb{R}^{m\times n}$ supports:
\begin{itemize}[leftmargin=*]
    \item Addition/subtraction with same shape
    \item Scalar product by $\lambda$
    \item Matrix-vector product:
    $$
    (Ax)_i = \sum_{j=0}^{n-1} a_{ij}x_j, \quad i=0,\dots,m-1
    $$
    \item Matrix-matrix product ($A\in\mathbb{R}^{m\times n}$, $B\in\mathbb{R}^{n\times p}$):
    $$
    (AB)_{ij} = \sum_{k=0}^{n-1} a_{ik}b_{kj}
    $$
\end{itemize}

Identity matrix $I_n$ is defined by:
$$
(I_n)_{ij}=
\begin{cases}
1, & i=j \\
0, & i\neq j
\end{cases}
$$

Suggested pseudocode for matrix-vector product:
\begin{verbatim}
for i = 0..m-1:
    sum = 0
    for j = 0..n-1:
        sum += A(i,j) * x[j]
    y[i] = sum
\end{verbatim}

\subsection*{b) Implementation}
Implement methods and operators declared in \texttt{LinAlgebra.hpp} and \texttt{LinAlgebra.cpp}:
\begin{itemize}[leftmargin=*]
    \item \texttt{Vector}: \texttt{operator+=}, \texttt{operator-=}, \texttt{operator*=}, \texttt{dot}
    \item Free operators for vector unary/binary operations and scalar operations
    \item \texttt{Matrix::identity}
    \item \texttt{Matrix::operator()(row,col)} read/write
    \item Matrix unary/binary operators, scalar operators, \texttt{Matrix*Vector}, \texttt{Matrix*Matrix}
\end{itemize}

\textbf{NOTE:} \textit{Always validate incompatible dimensions and invalid operations.}

\subsection*{c) Report}
\begin{enumerate}[leftmargin=*]
    \item Which dimension checks are required for each vector/matrix operator?
    \item Show one example where an invalid operation must throw an exception.
    \item Compare the computational cost of $Ax$ versus $AB$ using the
    operation counters (add, subtract, multiply, divide). Provide one concrete
    example for each case and explain, from the counter results, which operation
    grows faster when matrix size increases.
    \item Open test \texttt{LinearAlgebraTest.AlgebraicExpressionExpApprox}
    and inspect the implemented expression
    $$
    (I + A + \tfrac{1}{2}A^2)x + \tfrac{1}{4}(Bv).
    $$
    Report: (i) where this test is located, (ii) the numerical value obtained
    for the result vector with the provided data, and (iii) how you printed
    that value while running the test.
\end{enumerate}

\section*{Exercise 2: Elimination, Back Substitution, and Residual}

\subsection*{a) Mathematical Summary}
Given $Ax=b$ with $A\in\mathbb{R}^{n\times n}$, eliminate entries below each pivot.
At elimination step $k$:
$$
m_{ik} = \frac{a^{(k)}_{ik}}{a^{(k)}_{kk}},\quad i=k+1,\dots,n-1
$$
$$
a^{(k+1)}_{ij} = a^{(k)}_{ij} - m_{ik}a^{(k)}_{kj},\quad
i=k+1,\dots,n-1,\ j=k+1,\dots,n-1
$$
$$
b^{(k+1)}_i = b^{(k)}_i - m_{ik}b^{(k)}_k
$$

After elimination, solve $Ux=c$ by back substitution:
$$
x_i = \frac{c_i - \sum_{j=i+1}^{n-1} u_{ij}x_j}{u_{ii}},\quad i=n-1,\dots,0
$$

Residual quality metric:
$$
\text{relativeResidual}(A,x,b)=
\begin{cases}
\dfrac{\|Ax-b\|_2}{\|b\|_2}, & \|b\|_2 \neq 0 \\
\|Ax-b\|_2, & \|b\|_2 = 0
\end{cases}
$$

Suggested pseudocode (elimination core):
\begin{verbatim}
for k = 0..n-2:
    if |A(k,k)| <= tol: return false
    for i = k+1..n-1:
        m = A(i,k) / A(k,k)
        A(i,k) = m
        for j = k+1..n-1:
            A(i,j) -= m * A(k,j)
        b[i] -= m * b[k]
\end{verbatim}

\subsection*{b) Implementation}
Implement the next functions in \texttt{Solver.cpp}:
\begin{itemize}[leftmargin=*]
    \item \texttt{gaussianElimination(Matrix\& A, Vector\& b, bool pivoting, double tolerance, Operations\& operations)}
    \item \texttt{backSubstitution(const Matrix\& U, Vector\& c, double tolerance, Operations\& operations)}
    \item \texttt{relativeResidual(const Matrix\& A, const Vector\& x, const Vector\& b)}
\end{itemize}

For this exercise, focus on: validating sizes, pivots, multipliers, and row updates.

\subsection*{c) Report}
\begin{enumerate}[leftmargin=*]
    \item Why can a small pivot produce numerical instability?
    \item Why does back substitution ignore values below the diagonal in $U$?
    \item Explain why residual operations are not counted in solver operation counters.
    \item Provide one test case where elimination should return \texttt{false}.
\end{enumerate}

\section*{Exercise 3: Partial Pivoting}

\subsection*{a) Mathematical Summary}
At each elimination column $k$, choose:
$$
p = \arg\max_{i\in\{k,\dots,n-1\}} |a_{ik}|
$$
If $|a_{pk}| \leq \text{tolerance}$, stop with failure. Otherwise, if $p\neq k$,
swap rows $k$ and $p$ in both $A$ and $b$ before elimination updates.

Pivoting reduces sensitivity to round-off and avoids dividing by very small pivots
when a better row exists.

Suggested pseudocode:
\begin{verbatim}
for k = 0..n-2:
    p = argmax_i |A(i,k)| for i in [k, n-1]
    if |A(p,k)| <= tol: return false
    if p != k:
        swap row k and row p in A
        swap b[k] and b[p]
    continue with elimination step
\end{verbatim}

\subsection*{b) Implementation}
In \texttt{Solver.cpp}, extend \texttt{gaussianElimination}:
\begin{itemize}[leftmargin=*]
    \item If \texttt{pivoting == true}, choose and swap the pivot row.
    \item Preserve behavior with \texttt{pivoting == false}.
    \item Keep operation counters only for arithmetic operations, as specified.
\end{itemize}

\subsection*{c) Report}
\begin{enumerate}[leftmargin=*]
    \item Compare expected robustness with and without pivoting.
    \item Give one system where no-pivot fails but pivoting succeeds.
    \item Explain the additional computational overhead of pivot search.
    \item How does pivot tolerance interact with pivoting decisions?
\end{enumerate}

\section*{How to Test the Implementation}

Use the provided presets from the \texttt{lab1/} folder.

\subsection*{Assignment implementation}
\begin{verbatim}
cmake --preset default-tests
cmake --build --preset default-tests --parallel
ctest --preset default-tests
\end{verbatim}

Also test interactively with the UI. First, compile the executable:
\begin{verbatim}
cmake --preset default
cmake --build --preset default --parallel
\end{verbatim}

Then, run the app:
\begin{verbatim}
./build/bin/lab1
\end{verbatim}

Recommended testing strategy:
\begin{itemize}[leftmargin=*]
    \item Start with very small systems ($2\times2$, $3\times3$).
    \item Add edge cases: incompatible dimensions, near-singular pivots, and zero right-hand side.
    \item Verify exceptions and boolean failure paths.
    \item Compare residual quality against expected thresholds in datasets.
\end{itemize}

\section*{Evaluation Criteria}
Each exercise has a specific score:
\begin{itemize}[leftmargin=*]
    \item Exercise 1 (35\%): Implementation 25\% + Report 10\%
    \item Exercise 2 (35\%): Implementation 25\% + Report 10\%
    \item Exercise 3 (30\%): Implementation 20\% + Report 10\%
\end{itemize}

The Lab Assignment is optional. The final Lab grade is computed as the average between this
assignment submission and the Lab exam grade.
\begin{itemize}[leftmargin=*]
    \item If the student does not submit this assignment, the final Lab grade is the exam grade.
    \item If the student submits this assignment, the average between assignment and exam is used.
    However, if the average is lower than the Lab exam grade, the Lab exam grade is kept
    (so the student is not penalized).
\end{itemize}

Note: The code must work correctly. Otherwise, the grade of the corresponding
exercise is zero. If plagiarism is detected, the final grade is zero.

\section*{Academic Integrity and AI Use}
Using AI tools (LLMs, code assistants) as tutoring support is allowed.
Submitting generated code or report text as your own work without disclosure is
not allowed.

For submission:
\begin{itemize}[leftmargin=*]
    \item If AI support was used, include a short file \texttt{AI\_USAGE.md}
    in the submitted zip.
    \item In \texttt{AI\_USAGE.md}, specify: tool name, what it was used for,
    and the final prompts that directly influenced the submitted code or report.
    \item You must be able to explain and justify every submitted function and
    test.
\end{itemize}

During review, the student will be asked for a brief explanation of selected code fragments or test decisions. If a submitted part cannot be explained by the
student, that part will be considered invalid for grading.

\section*{Submission}
For submission, you may form groups of up to 3 people. Through ATENEA,
use the format:
\begin{center}
\texttt{Surname1Name1\_Surname2Name2\_Surname3Name3\_Lab1.zip}
\end{center}
with the following content:
\begin{itemize}[leftmargin=*]
    \item Code files:
    \begin{itemize}[leftmargin=*]
        \item \texttt{LinAlgebra.cpp} with the implementation of classes \texttt{Vector} and \texttt{Matrix}.
        \item \texttt{Solver.cpp} with the implementation of the solving methods \texttt{Solver}.
    \end{itemize}
    \item Report (PDF, 2--4 pages).
\end{itemize}

The submission will be evaluated jointly with the group members.
The lab exam will be individual.

\section*{Important Dates}
% \vspace{2\baselineskip}
\noindent\fbox{%
\parbox{\dimexpr\linewidth-2\fboxsep-2\fboxrule\relax}{%
\centering\Large\bfseries
Submission Deadline: October 27, 2026.\\
Lab Practical Exam: October 29, 2026.
}%
}
\end{document}
